% Reconstrucción del enunciado de la Tarea 1 (a partir de assignment-1.pdf)
% Semestre actualizado a 2026-2 y fecha de compilación (\today).
\documentclass[11pt,spanish]{article}

\input{preamble.tex}

% --- Datos del enunciado ---
\newcommand{\tnum}{1}
\newcommand{\sem}{2026-2}
\newcommand{\campus}{Casa Central \\ Valparaíso}
\newcommand{\deadline}{8 de Septiembre de 2026}

% --- Encabezado y pie de página ---
\headheight=14pt
\linespread{1.3}
\pagestyle{fancy}
\fancyhf{}
\lhead{INF-221 TAREA \tnum}
\rhead{\sem}
\fancyfoot[L]{Campus \campus}
\fancyfoot[C]{\thepage}
\fancyfoot[R]{ALGORITMOS Y COMPLEJIDAD \sem}
\renewcommand{\headrulewidth}{0.4pt}
\renewcommand{\footrulewidth}{0.4pt}

\title{%
  \huge
  \textbf{ENUNCIADO TAREA \tnum~ \\ ALGORITMOS Y COMPLEJIDAD} \\[1ex]
  \emph{\textquote{Más allá de la notación asintótica: Análisis experimental de algoritmos de ordenamiento y multiplicación de matrices.}}
}

\author{%
  Fecha límite de entrega: \deadline \\[1ex]
  Equipo ALGORITMOS Y COMPLEJIDAD \sem
}

\date{%
  \small
  \today \\
  \currenttime \\[1ex]
  Versión 0.0
}

\begin{document}
\maketitle
\thispagestyle{fancy}
\vspace{-1.0\baselineskip}

\newpage

\tableofcontents

\newpage

\section{Objetivos}
El objetivo de esta tarea es introducirnos en el Análisis experimental de Algoritmos. Para ello, se propone realizar un estudio experimental de 4 algoritmos de ordenamiento y 2 de multiplicación de matrices.

Para cada uno de los dos problemas, el de ordenar un arreglo unidimensional de números enteros, y el de multiplicar dos matrices cuadradas, existen múltiples algoritmos que los resuelven, y la mayoría, si no todos, tienen implementaciones de fácil acceso (ya se encuentran en bibliotecas de lenguajes de programación, o son fáciles de encontrar en Internet). Sus complejidades teóricas son conocidas, ya que estás son una característica fundamental a la hora de diseñar algoritmos y, en la mayoría de los casos, la motivación para diseñarlos.

El desafío en esta tarea no está en diseñar los algoritmos (los encuentra fácilmente), ni tampoco en calcular o demostrar su complejidad teórica: el desafío está en \textbf{estudiar su comportamiento en la práctica} y relacionar este comportamiento con la complejidad teórica.

\section{Especificaciones}
En esta sección se describen los pasos a seguir para completar la tarea, la cual consiste en desarrollar código en C++ \hyperref[subsec:implementaciones]{«Implementaciones»} y en realizar un mini-informe \hyperref[subsec:mini-informe]{«Mini-informe»}. Se espera que cada uno de los pasos se realice de manera ordenada y siguiendo las instrucciones dadas.

\subsection{Implementaciones}\label{subsec:implementaciones}

\begin{mdframed}[linecolor=blue!60!black, backgroundcolor=blue!5,
                 linewidth=1pt, roundcorner=5pt, innertopmargin=8pt,
                 innerbottommargin=8pt]
\textbf{Nota: }Abra este documento en algún lector de pdf que permita hipervínculos, ya que en este documento el texto en color azul suele indicar un hipervínculo.
\end{mdframed}


\begin{enumerate}[(1)]
  \item En caso de dudas, enviar preguntas \textbf{al foro de la Tarea 1}. En caso de cualquier de modificaciones, todas se informarán por aula.
  \item Todo lo necesario para realizar la tarea se encuentra en el archivo \verb|.zip| proporcionado a través de aula.usm.cl.
\end{enumerate}

\begin{mdframed}[linecolor=blue!60!black, backgroundcolor=blue!5,
                 linewidth=1pt, roundcorner=5pt, innertopmargin=8pt,
                 innerbottommargin=8pt]
                 \textbf{Nota:} La estructura de archivos y el template del proyecto se encuentran disponibles en el material entregado.
\end{mdframed}


\newpage

\begin{enumerate}[(1)]
  \item \textbf{Ordenamiento y Multiplicación.}

  \textbf{Ordenamiento.} Implementar en C++ los algoritmos de ordenamiento: \textsc{merge sort}, \textsc{quick sort}, \textsc{patience sort} y el algoritmo \textsc{sort} de la librería estándar de C++. Use la siguiente estructura de archivos:
  \begin{itemize}
    \item \verb|code/sorting/algorithms/mergesort.cpp|
    \item \verb|code/sorting/algorithms/quicksort.cpp|
    \item \verb|code/sorting/algorithms/patiencesort.cpp|
    \item \verb|code/sorting/algorithms/sort.cpp|
  \end{itemize}
  El algoritmo \textsc{sort} ya se encuentra implementado en el material entregado con el fin de que sea incluido en sus mediciones del punto (2) de la subsección \hyperref[subsec:implementaciones]{«Implementaciones»}.

  \textbf{Multiplicación.} Se deben implementar en C++ los algoritmos de multiplicación \textsc{naive} y \textsc{strassen}. Use la siguiente estructura de archivos.
  \begin{itemize}
    \item \verb|code/matrix_multiplication/algorithms/naive.cpp|
    \item \verb|code/matrix_multiplication/algorithms/strassen.cpp|
  \end{itemize}

  \item \textbf{Mediciones de tiempo y memoria.} Implementar el programa que realizará las mediciones de tiempo y memoria en C++ (programas principales) y su respectivo makefile, que ejecutará los algoritmos y generará los archivos de salida en cada uno de los directorios measurements sorting y measurements matrix multiplication con los resultados de cada uno de los algoritmos.
  \begin{itemize}
    \item \verb|code/sorting/sorting.cpp|
    \item \verb|code/matrix_multiplication/matrix_multiplication.cpp|
  \end{itemize}

  \item \textbf{Generación de gráficos.} Implementar el programa que generará los gráficos en \textsc{python} y que se encargará de leer los archivos generados por los programas principales guardados en measurements sorting y measurements matrix multiplication, para luego graficar los resultados obtenidos y guardarlos en formato PNG en plots sorting y plots matrix multiplication.
  \begin{itemize}
    \item \verb|code/sorting/scripts/plot_generator.py|
    \item \verb|code/matrix_multiplication/scripts/plot_generator.py|
  \end{itemize}

  \item \textbf{Documentación.} Documentar cada uno de los pasos anteriores
  \begin{itemize}
    \item Completar el archivo README.md del directorio code
    \item Documentar en cada uno de sus programas, al inicio de cada archivo, fuentes de información, referencias y/o bibliografía utilizada para la implementación de cada uno de los algoritmos.
  \end{itemize}
\end{enumerate}

\newpage

\subsection{Mini-informe}\label{subsec:mini-informe}
Generar un mini-informe en \LaTeX{} que contenga una introducción breve, los resultados obtenidos (tablas y gráficos) con una discusión sobre ellos, y una conclusión breve. En el material entregado podrá encontrar el Template en el directorio report que deben utilizar:

\begin{center}
  \verb|Template disponible en: report/ dentro del material entregado|
\end{center}

Las indicaciones se encuentran en el archivo README.md y en la plantilla de \LaTeX{}.

\section{Condiciones de entrega}
\begin{enumerate}[(1)]
  \item La tarea es individual, sin excepciones.
  \item La entrega debe realizarse vía aula.usm.cl, subiendo un archivo \verb|.zip| que contenga todo el contenido de su proyecto (código, informe, datos, gráficos, etc.).
  \item Si se utiliza algún código, idea, o contenido extraído de otra fuente, este debe ser citado en el lugar exacto donde se utilice.
  \item Asegúrese que todas sus entregas tengan sus datos completos: número de la tarea, ramo, semestre, nombre y rol. Puede incluirlas como comentarios en sus fuentes \LaTeX{} (en \TeX{} comentarios son desde \% hasta el final de la línea) o en posibles programas. Anótese como autor de los textos.
  \item Si usa material adicional al discutido en clases, detállelo. Agregue información suficiente para ubicar ese material (en caso de no tratarse de discusiones con compañeros de curso u otras personas).
  \item La fecha límite de entrega es el día 9 de Abril de 2026. \\ \textbf{NO SE ACEPTARÁN TAREAS FUERA DE PLAZO.}
  \item Nos reservamos el derecho de llamar a interrogación sobre algunas de las tareas entregadas. En tal caso, la nota de la tarea será la obtenida en la interrogación. \\ \textbf{NO PRESENTARSE A UN LLAMADO A INTERROGACIÓN SIN JUSTIFICACIÓN PREVIA SIGNIFICA AUTOMÁTICAMENTE NOTA 0.}
\end{enumerate}

\newpage

\appendix

\section{Casos de prueba}

\subsection{Ordenamiento de un arreglo unidimensional de números enteros}

\textbf{Entrada:} Leer un arreglo unidimensional $A$ desde el archivo \verb|{n}_{t}_{d}_{m}.txt|.

\begin{itemize}
  \item $n$ hace referencia a la cantidad de elementos (o largo del arreglo) y pertenece al conjunto $N = \{10^1, 10^3, 10^5, 10^7\}$.
  \item $t$ hace referencia al tipo de matriz, y pertenece al conjunto $T = \{\text{ascendente}, \text{descendente}, \text{aleatorio}\}$.
  \item $d$ hace referencia al conjunto dominio de cada elemento del arreglo. $d = \{D1, D7\}$, donde $D7$ implica que el dominio es $\{0,1,2,3,4,5,6,7,8,9\}$ y $D7$ que el dominio es $\{0,1,2,3,\dots,10^7\}$.
  \item $m$ hace referencia a la muestra aleatoria (o caso de prueba) y pertenece al conjunto $M = \{a, b, c\}$.
\end{itemize}

\textbf{Salida:} Escribir la matriz resultante $M_1 \times M_2$ en el archivo \verb|{n}_{t}_{d}_{m}_out.txt|.

Los programas que generan estos arreglos se encuentran en

\begin{center}
  \verb|code/sorting/scripts/array_generator.py|
\end{center}

\subsection{Multiplicación de matrices cuadradas de números eneteros}

\textbf{Entrada:} Leer dos matrices cuadradas de entrada $M_1$ y $M_2$ desde los archivos \verb|{n}_{t}_{d}_{m}_1.txt| y \verb|{n}_{t}_{d}_{m}_2.txt|, respectivamente.

\begin{itemize}
  \item $n$ hace referencia a la dimensión de la matriz ($n$ filas y $n$ columnas) y pertenece al conjunto $N = \{2^4, 2^6, 2^8, 2^{10}\}$.
  \item $t$ hace referencia al tipo de matriz, y pertenece al conjunto $T = \{\text{dispersa}, \text{diagonal}, \text{densa}\}$.
  \item $d$ hace referencia al dominio de cada coeficiente de la matriz $d = \{D0, D10\}$, donde $D0$ implica que el dominio es $\{0,1\}$ y $D10$ que el dominio es $\{0,1,2,3,4,5,6,7,8,9\}$.
  \item $m$ hace referencia a la muestra aleatoria (o caso de prueba) y pertenece al conjunto $M = \{a, b, c\}$.
\end{itemize}

\textbf{Salida:} Escribir la matriz resultante $M_1 \times M_2$ en el archivo \verb|{n}_{t}_{d}_{m}_out.txt|.

Los programas que generan estas matrices se encuentran en

\begin{center}
  \verb|code/matrix_multiplication/scripts/matrix_generator.py|
\end{center}

\end{document}
